\documentclass[10pt,a4paper]{amsart}
\usepackage[margin=19mm]{geometry}
\usepackage{amsmath,amssymb,mathtools}
\usepackage[scr=rsfs,cal=euler]{mathalfa}
\usepackage{xfrac}
\usepackage[unicode,hidelinks,bookmarks=false]{hyperref}
\newcommand{\ie}{i.e.\ }
\newcommand{\cf}{cf.\ }
\newcommand{\resp}{respectively\ }
\newcommand{\Subsection}{\subsection}
\providecommand{\nobreakdash}{\nobreak\hbox{-}}
\setlength{\emergencystretch}{2em}
\multlinegap0pt
\usepackage{bm}
\usepackage{bbm}
\usepackage{dsfont}
\usepackage{xspace}
\usepackage{cancel}
\usepackage{mathtools}
\usepackage{amsthm}
\makeatletter
\@ifundefined{theo}{\newtheorem{theo}{Theo}}{}
\@ifundefined{defn}{\newtheorem{defn}[theo]{Definition}}{}
\@ifundefined{prop}{\newtheorem{prop}[theo]{Proposition}}{}
\makeatother
\def \ie{{\it i.e.}}
\def\QQ{{\mc Q}}
\newcommand{\Alg}{{\textup{Alg}}}
\newcommand{\CC}{{\textup{CC}}}
\newcommand{\Def}{{\textup{Def}}}
\newcommand{\Kh}{{\K[\hbar]}}
\newcommand{\K}{{\bb K}}
\newcommand{\MC}{{{\textup{MC}}}}
\newcommand{\Obs}{{\textup{Obs}}}
\newcommand{\Ph}{{\mathcal{P}_\hbar}}
\newcommand{\Qh}{{\mathcal{Q}_\hbar}}
\newcommand{\Spec}{{\textup{Spec\ }}}
\newcommand{\Ss}{{{\mf{S}}}}
\newcommand{\ash}{{\textup{\mbox{!`}}}}
\newcommand{\bb}{\mathbb}
\newcommand{\fDef}{{\mathfrak{Def}}}
\newcommand{\fMC}{{{\mathfrak{MC}}}}
\newcommand{\fibre}[1]{\underset{#1}{\times}}
\newcommand{\gPQh}{{\mathfrak{g}_{\mc P_\hbar,\mc Q_\hbar}}}
\newcommand{\g}{{\mf{g}}}
\newcommand{\ik}{{\bb K}}
\newcommand{\mc}{\mathcal}
\newcommand{\mf}{\mathfrak}
\title[On the deformation theory of chiral quantizations]{On the deformation theory of chiral quantizations}
\author{Dylan Butson, Sujay Nair}
\date{Source: arXiv:2606.27341v1 (CC BY 4.0)}
\begin{document}
\maketitle

\noindent\emph{Source and licence.} On the deformation theory of chiral quantizations. Version arXiv:2606.27341v1 (CC BY 4.0), distributed under the Creative Commons Attribution 4.0 International licence, as recorded in the arXiv licence metadata for this exact version and shown on the abstract page https://arxiv.org/abs/2606.27341v1. Full rights notice: licenses/arxiv-2606.27341-CC-BY-4.0.txt.\par\smallskip

\noindent\textit{Editorial scope.} Counted unit: the Theo whose source label is \texttt{DefObsThm} in the article (source0.tex lines 3626--3635, source label \texttt{DefObsThm}), delivered together with the article's own complete original proof of it (source0.tex lines 3637--3709, one \texttt{proof} environment of 73 lines). No mathematics is written, repaired or re-derived: statement and proof are the article's own text line for line. Required context retained uncounted: MCstackeqn (lines 3382--3384); DefObsprop (lines 3431--3449); Defhbardefn (lines 3553--3567); Ssmodtrivdefn (lines 3597--3603); gtriveqn (lines 3606--3608); Obshbareqn (lines 3621--3623). Every definition, display and label of those lines is retained unchanged. Omitted from this package and not redistributed: 1--3381, 3385--3430, 3450--3552, 3568--3596, 3604--3605, 3609--3620, 3624--3625, 3636--3636, 3710--6110. Nothing inside a retained excerpt is omitted or altered: every retained body is delivered exactly as the article's lines are written, so no omission receipt of any kind accompanies this package. Citations: the retained excerpts carry no citation command at all, so no bibliography entry of the article is transcribed and no reference is redistributed. Reference adaptation: none was needed and none was made; every reference inside the delivered unit resolves inside it, so no unresolved reference or citation is produced. Printed slips kept literal: no printed word, symbol, hypothesis or number of the counted statement or of its proof was altered. Layout and class: the article's own class is \texttt{amsart} with packages including amsmath, amsthm, amssymb, hyperref, xy, xcolor, mleftright, overpic, geometry, wrapfig, bbm, arydshln, multirow, mathtools; the wrapper is the lane's pinned \texttt{amsart} editorial wrapper at 10pt on A4, which restates the theorem-like environments the retained text needs and adds the article's own macro definitions that the retained text needs (\texttt{\textbackslash Alg}, \texttt{\textbackslash CC}, \texttt{\textbackslash Def}, \texttt{\textbackslash K}, \texttt{\textbackslash Kh}, \texttt{\textbackslash MC}, \texttt{\textbackslash Obs}, \texttt{\textbackslash Ph}, \texttt{\textbackslash QQ}, \texttt{\textbackslash Qh}, \texttt{\textbackslash Spec}, \texttt{\textbackslash Ss}, \texttt{\textbackslash ash}, \texttt{\textbackslash bb}, \texttt{\textbackslash fDef}, \texttt{\textbackslash fMC}, \texttt{\textbackslash fibre}, \texttt{\textbackslash g}, \texttt{\textbackslash gPQh}, \texttt{\textbackslash ie}, \texttt{\textbackslash ik}, \texttt{\textbackslash mc}, \texttt{\textbackslash mf}), with extra wrapper packages (bm, bbm, dsfont, xspace, cancel, mathtools). The delivered numbering is the unit's own. Assets: no retained span contains a figure environment, a graphics-inclusion command, a PDF-page inclusion, a table or a TikZ picture; no image file is copied into this package, the package declares no asset dependency and no image file is redistributed with it. Licence: version arXiv:2606.27341v1 of this article is distributed under the Creative Commons Attribution 4.0 International licence, as recorded in the arXiv licence metadata for this exact version and shown on the abstract page https://arxiv.org/abs/2606.27341v1. A full rights notice is delivered at licenses/arxiv-2606.27341-CC-BY-4.0.txt. Characters: every retained excerpt is pure ASCII, so the delivered engine, XeLaTeX with its default Latin Modern text font, reports no missing character.
\begin{equation}\label{MCstackeqn}
		\fMC_\varphi(S) =\left[ \MC( \g^\varphi_{\mc P, \mc Q}\otimes_{\bb K} \mf m_S)^{(1)}/\mf G_{\mc P,\mc Q,\mf m_S}^{0,(0)}\right]   \ .
	\end{equation}

\begin{prop}\label{DefObsprop} 
Let $\varphi \in \Alg_{\mc P}(\mc Q)$ and $\varphi_{\hbar/\hbar^{m}}\in \Def_\varphi( \K[\hbar]/\hbar^{m})$. Then the obstruction associated to $\varphi_{\hbar/\hbar^m}$ is closed, \ie,
\[
d_\varphi \Obs_m(\varphi_{\hbar/\hbar^m}) =0 
\]
Moreover, for some putative extension to $m$th order in $\hbar$, $\varphi_{\hbar/\hbar^{m+1}} = \varphi_{\hbar/\hbar^m}+ \hbar^m \varphi_m$, we have 
\[ 
\varphi_{ / \hbar^{m+1}} \in \Def_\varphi(\K[\hbar]/\hbar^{m+1}) \times_{\Def_\varphi(\K[\hbar]/\hbar^{m}) } \{\varphi_{\hbar/\hbar^{m}} \} \qquad\text{if and only if} \qquad d_\varphi \varphi_m = \Obs_m(\varphi_{\hbar/\hbar^{m}})  \ .
\]


Thus, there exists an extension
\[ \varphi_{ \hbar/ \hbar^{m+1}} \in \Def_\varphi(\K[\hbar]/\hbar^{m+1}) \times_{\Def_\varphi(\K[\hbar]/\hbar^{m}) } \{\varphi_{\hbar/\hbar^{m}} \} \]
if and only if the cohomology class of the obstruction vanishes, that is,
\[ [ \Obs_m(\varphi_{\hbar/\hbar^{m}})]= 0 \ \in H^2_{\mc P,\mc Q}(\varphi)\ .\]
Further, when the obstruction class vanishes, the space of such extensions up to isomorphism,
\[  \fDef_\varphi(\K[\hbar]/\hbar^{m+1}) \times_{\fDef_\varphi(\K[\hbar]/\hbar^{m}) } \{\varphi_{\hbar/\hbar^{m}} \} \ , \]
is a torsor for $H^1_{\mc P,\mc Q}(\varphi)$.
\end{prop}

\begin{defn}\label{Defhbardefn} 
Let $\varphi\in \Alg_{\mc P_0}(\mc Q_0)$ be an $\mc P$-algebra internal to $\mc Q$. The space of deformations of $\varphi$ over $\Spec S$, as an $\mc P_\hbar$-algebra internal to $\mc Q_\hbar$, is defined as the fibre product
\[ 
\Def_\varphi^{\Ph,\Qh}(S) = \MC(\gPQh\otimes_{\K[\hbar]} S) \fibre{\MC(\g_{\mc P_0,\mc Q_0})} \{\varphi\} \ .  
\]
Similarly, the stack of such deformations is defined as the fibre product of Maurer-Cartan stacks
\[
\fDef_{\varphi}^{\Ph,\Qh}(S) = \fMC(\gPQh\otimes_{\K[\hbar]} S) \fibre{\fMC(\g_{\mc P_0,\mc Q_0})} \{\varphi_0\} \ , 
\]
where the Maurer-Cartan stack is defined as the quotient stack
\[ 
\fMC(\g_{\Ph,\Qh }\otimes_{\K[\hbar]} S)= \left[ \MC(\g_{\Ph,\Qh }\otimes_{\K[\hbar]} S)\  \big/\  \mf G^{0,(0)}_{\Ph,\Qh,\mf m_S} \right]  
\]
and the group $\mf G^{0,(0)}_{\Ph,\Qh,\mf m_S}=\exp(\g^{0,(0)}_{\Ph,\Qh}\otimes_{\K[\hbar]} \mf m_S)$, the formal exponential of the nilpotent Lie algebra $\g^{0,(0)}_{\Ph,\Qh}\otimes_{\K[\hbar]} \mf m_S$---analogous to the definition of the Maurer-Cartan stack in Equation \ref{MCstackeqn}.
\end{defn}

\begin{defn}\label{Ssmodtrivdefn} 
By a trivialization of the underlying $\Ss$-module of an operad $\QQ_\hbar$ or a cooperad $\CC_\hbar$ over $\ik[\hbar]$ we mean a choice of trivialization at each arity $n$ of the $\ik[\hbar]$-modules $\QQ_\hbar(n)$ and $\CC_\hbar(n)$, \ie am isomorphism of $\ik[\hbar]$-modules
\[
\QQ_{\hbar}^\ash(n) \cong \QQ_0^\ash(n) \otimes \ik[\hbar] \qquad \text{ and } \qquad \CC_\hbar(n) \cong \CC_0(n) \otimes \ik[\hbar]~.
\]
Note that this is much weaker than a trivialization of the operadic or cooperadic structure.
\end{defn}

\begin{equation}\label{gtriveqn}
	 \gPQh \cong \g_{\mc P_0,\mc Q_0}\otimes_{\K} \K[\hbar] \qquad\text{with}\qquad d_{\gPQh}= \sum_{j=0}^\infty \hbar^j d_j \quad \text{and} \quad [\cdot ,\cdot]_\gPQh = \sum_{j=0}^\infty \hbar^j[\cdot,\cdot]_j \ ,
\end{equation}

\begin{equation}\label{Obshbareqn}
	 \Obs_m^{\Ph,\Qh}(\varphi_{\hbar/\hbar^{m}}):=-\frac{1}{2} \sum_{k=1}^{m-1}[ \varphi_k,\varphi_{m-k} ] -  \sum_{j=1}^m d_j \varphi_{m-j} - \sum_{j=1}^m \sum_{k=0}^{m-j} \frac{1}{2}[\varphi_k, \varphi_{m-j-k}]_j  \ \in \g^{2,(2)}_{\mc P,\mc Q} \ .
\end{equation}

\begin{theo}\label{DefObsThm} Let $\varphi\in \Alg_{\mc P_0}(\mc Q_0)$ be a $\mc P_0$-algebra internal to $\mc Q_0$ and $\varphi_{\hbar/\hbar^{m}}\in \Def^{\Ph,\Qh}_\varphi(\K[\hbar]/\hbar^{m}) $. Then $\Obs_m(\varphi_{\hbar/\hbar^{m}})$ satisfies $d_\varphi \Obs_m(\varphi_{\hbar/\hbar^{m}})=0$, and for  $\varphi_{ \hbar/ \hbar^{m+1}} =  \varphi_{\hbar/\hbar^{m}} + \hbar^m \varphi_m $ we have
		\[ \varphi_{ \hbar/ \hbar^{m+1}} \in \Def^{\Ph,\Qh}_\varphi(\K[\hbar]/\hbar^{m+1}) \times_{\Def^{\Ph,\Qh}_\varphi(\K[\hbar]/\hbar^{m}) } \{\varphi_{\hbar/\hbar^{m}} \} \qquad \text{if and only if} \qquad d_{\varphi} \varphi_m  = \Obs_m(\varphi_{\hbar/\hbar^{m}}) \ . \]
		Thus, there exists an extension
\[  \varphi_{ \hbar/ \hbar^{m+1}} \in \Def^{\Ph,\Qh}_\varphi(\K[\hbar]/\hbar^{m+1}) \times_{\Def^{\Ph,\Qh}_\varphi(\K[\hbar]/\hbar^{m}) } \{\varphi_{\hbar/\hbar^{m}} \} \]
if and only if the cohomology class of the obstruction vanishes, that is,
\[ [ \Obs^{\Ph,\Qh}_m(\varphi_{\hbar/\hbar^{m}})]= 0 \ \in H^{2}_{\mc P_0,\mc Q_0}(\varphi) \ .\]
Further, when the obstruction class vanishes, the space of such extensions up to isomorphism,
\[  \fDef^{\Ph,\Qh}_\varphi(\K[\hbar]/\hbar^{m+1}) \times_{\fDef^{\Ph,\Qh}_\varphi(\K[\hbar]/\hbar^{m}) } \{\varphi_{\hbar/\hbar^{m}} \} \ , \]
is a torsor for $H^{1}_{\mc P_0,\mc Q_0}(\varphi)$.
\end{theo}

\begin{proof}
We repeat the argument of Proposition \ref{DefObsprop}, keeping track of the additional $\hbar$-dependence coming from the families $\Ph$ and $\Qh$, and the fixed trivializations of their underlying $\Ss$-modules in the sense of Definition \ref{Ssmodtrivdefn}. In particular, via these trivializations one has an identification of $\Kh$-modules
\[
\gPQh \;\cong\; \g_{\mc P_0,\mc Q_0}\otimes_\K \Kh
\]
under which the differential and bracket on $\gPQh$ admit expansions
\[
d_{\gPQh}=\sum_{j\ge 0}\hbar^j d_j
\qquad\text{and}\qquad
[\cdot,\cdot]_{\gPQh}=\sum_{j\ge 0}\hbar^j[\cdot,\cdot]_j
\]
for $d_j$ and $[\cdot,\cdot]_j$ defined on the underlying vector space $\g_{\mc P_0,\mc Q_0}$, as in the description following Equation \ref{gtriveqn}.

Fix $m\ge 1$ and suppose that
\[
\varphi_{\hbar/\hbar^m}=\varphi+\sum_{k=1}^{m-1}\hbar^k\varphi_k
\;\in\;(\gPQh)^{1,(1)}\otimes_{\Kh}\Kh/(\hbar^m)
\]
is a Maurer-Cartan element, that is
\[
d_{\gPQh}(\varphi_{\hbar/\hbar^m})
+\frac12[\varphi_{\hbar/\hbar^m},\varphi_{\hbar/\hbar^m}]_{\gPQh}= 0 \ \in (\gPQh)^{2,(2)}\otimes_{\Kh}\Kh/(\hbar^{m}).
\]
Consider an arbitrary lift
\[
\varphi_{\hbar/\hbar^{m+1}}
=\varphi+\sum_{k=1}^{m}\hbar^k\varphi_k
\ \in (\gPQh)^{1,(1)}\otimes_{\Kh}\Kh/(\hbar^{m+1}),
\]
and expand the Maurer-Cartan equation modulo $\hbar^{m+1}$. The Maurer-Cartan equations for $\varphi_{\hbar/\hbar^{m+1}}$ modulo $\hbar^m$ are precisely the same as those for $\varphi_{\hbar/\hbar^{m}}$, which hold by hypothesis. So it remains to check the vanishing of the coefficient of $\hbar^m$. Using the expansions of $d_{\gPQh}$ and $[\cdot,\cdot]_{\gPQh}$ fixed above and collecting the terms of order $\hbar^m$ gives
\[
d_0\varphi_m+[\varphi,\varphi_m]_0
+\frac12\sum_{k=1}^{m-1}[\varphi_k,\varphi_{m-k}]_0
+\sum_{j=1}^m d_j\varphi_{m-j}
+\sum_{j=1}^m\sum_{k=0}^{m-j}\frac12[\varphi_k,\varphi_{m-j-k}]_j
=0.
\]
By the Definition of $\Obs_m^{\Ph,\Qh}(\varphi_{\hbar/\hbar^m})$ in Equation \ref{Obshbareqn}, this is equivalent to
\begin{equation}\label{defMChbareqn}
	d_\varphi(\varphi_m)=\Obs_m^{\Ph,\Qh}(\varphi_{\hbar/\hbar^m}),
\end{equation}
where $d_\varphi:=d_0+[\varphi,-]_0$ is the differential on the operadic deformation complex
$C^\bullet_{\mc P_0,\mc Q_0}(\varphi)$. Thus, a lift $\varphi_{\hbar/\hbar^{m+1}}=\varphi_{\hbar/\hbar^m}+\hbar^m\varphi_m$ is Maurer-Cartan
modulo $\hbar^{m+1}$ if and only if $d_\varphi\varphi_m=\Obs_m^{\Ph,\Qh}(\varphi_{\hbar/\hbar^m})$, as claimed. In particular, an extension exists if and only if the cochain
$\Obs_m^{\Ph,\Qh}(\varphi_{\hbar/\hbar^m})$ is exact, equivalently
\[
[\Obs_m^{\Ph,\Qh}(\varphi_{\hbar/\hbar^m})]=0\in H^2_{\mc P_0,\mc Q_0}(\varphi).
\]



As in Proposition \ref{DefObsprop}, we use the Bianchi identity to show that $\Obs_m^{\Ph,\Qh}(\varphi_{\hbar/\hbar^m})$ is in general a $d_\varphi$-cocycle. Since $\varphi_{\hbar/\hbar^m}$ satisfies the Maurer-Cartan equation up to order $\hbar^m$, we have the expansion 
\begin{equation}\label{MChbarmpoeqn}
	d_\hbar\varphi_{\hbar/\hbar^{m+1}} \;+\; \frac{1}{2}[\varphi_{\hbar/\hbar^{m+1}},\varphi_{\hbar/\hbar^{m+1}}] \;=\; \hbar^{\,m}\,\Obs_m^{\Ph,\Qh}(\varphi_{\hbar/\hbar^m}) \;+\; \textup{terms of order } \hbar^{\,m+1}\,.
\end{equation}
The Bianchi identity modulo $\hbar^{m+1}$ is given by
\[ 
d_{\varphi_{ \hbar / \hbar^{m+1}}}\!\Big(d_\hbar\varphi_{\hbar/\hbar^{m+1}} \;+\; \frac{1}{2}[\varphi_{\hbar/\hbar^{m+1}},\varphi_{\hbar/\hbar^{m+1}}] \Big) \;=\; 0\,.
\] 
so that applying the differential $d_{\varphi_{ \hbar / \hbar^{m+1}}}$ to both sides of Equation \ref{MChbarmpoeqn}, the vanishing at order $\hbar^m$ is equivalent to
\[ 
d_\varphi\!\Big(\Obs_m^{\Ph,\Qh}(\varphi_{\hbar/\hbar^m})\Big) \;=\; 0\,,
\] 
and thus $\Obs_m^{\Ph,\Qh}(\varphi_{\hbar/\hbar^m})$ is indeed closed with respect to the differential $d_\varphi$.


Finally, suppose $[\Obs_m^{\Ph,\Qh}(\varphi_{\hbar/\hbar^m})]=0$ so that extensions exist. If $\varphi_{\hbar/\hbar^{m+1}}=\varphi_{\hbar/\hbar^m}+\hbar^m\varphi_m$ and $\tilde\varphi_{\hbar/\hbar^{m+1}}=\varphi_{\hbar/\hbar^m}+\hbar^m\tilde \varphi_m$ are two such extensions, then we have
\[
d_\varphi(\tilde\varphi_m-\varphi_m)=0,
\]
so $\delta:=\tilde\varphi_m-\varphi_m$ is a $1$-cocycle. Moreover, $\delta$ is a coboundary, $\delta=d_\varphi(\eta)$, if and only if the two lifts are related by the action of
$\mf G^{0,(0)}_{\Ph,\Qh,\mf m}$ at order $\hbar^m$, where the group $\mf G^{0,(0)}_{\Ph,\Qh,\mf m_R}$ is as in Definition \ref{Defhbardefn}. Thus, isomorphism classes of such extensions again form a torsor under $H^1_{\mc P_0,\mc Q_0}(\varphi)$, as claimed.
\end{proof}

\end{document}
